\chapter{Oligopoly and Cournot Competition}
\label{chap:oligopoly}

\section{Learning objectives}

After completing this chapter, you should be able to:
\begin{itemize}
  \item identify strategic interdependence in oligopoly;
  \item derive and interpret Cournot reaction functions;
  \item solve a two-firm Cournot--Nash equilibrium; and
  \item compare monopoly, oligopoly, and competition.
\end{itemize}

\section{Oligopoly and strategic interaction}

An \keyterm{oligopoly} is a market with a small number of firms. Each firm's
decision affects market outcomes and therefore the profits of its rivals. A
small set of mobile-network providers or airlines on a route can illustrate
the idea: one firm's price, capacity, or advertising decision changes the
environment facing the others. This \keyterm{strategic interdependence}
distinguishes oligopoly from both perfect competition and monopoly.

Firms may compete in quantities, prices, capacity, product characteristics,
advertising, or research. The appropriate model depends on the strategic
variable and timing. Cournot competition studies firms that choose output
quantities simultaneously. It is most useful when output or capacity is chosen
before sales occur, such as when firms commit production before customers
arrive. If firms instead choose prices directly, a \keyterm{Bertrand
competition} model may be more appropriate.

\begin{exercisebox}
\textbf{Two firms choosing capacity.}
Imagine two firms that must decide how much bottled water to prepare before
the market price is known. More total output lowers the eventual market price.
Each firm therefore wants to produce more when its rival produces less. This
is the logic captured by a downward-sloping reaction function.
\end{exercisebox}

\section{The Cournot model}

Suppose two firms produce a homogeneous good. Firm $i$ chooses $q_i$, total
output is $Q=q_1+q_2$, and inverse demand is $p(Q)$. When firm $i$ chooses
output, it treats the rival's current output $q_j$ as fixed. In effect, it
asks: given what the rival is producing now, what quantity is best for me?
This does not mean that the firm believes its rival will never react; each
firm solves a unilateral choice problem, and equilibrium requires the two
resulting best responses to be mutually consistent.

With
\[
p=80-2Q
\qquad\text{and}\qquad
C(q_i)=10+4q_i,
\]
firm $i$'s profit is
\[
\pi_i
=(80-2(q_i+q_j))q_i-10-4q_i
=76q_i-2q_i^2-2q_iq_j-10.
\]
Holding $q_j$ fixed,
\[
\frac{\partial\pi_i}{\partial q_i}
=76-4q_i-2q_j,
\qquad
\frac{\partial^2\pi_i}{\partial q_i^2}=-4<0.
\]
The first-order condition therefore gives the unique unconstrained
\keyterm{best response}:
\[
q_i=19-\frac{1}{2}q_j.
\]
Because output cannot be negative, the full \keyterm{reaction function} is
\[
BR_i(q_j)=\max\left\{0,\,
19-\frac{1}{2}q_j\right\}.
\]
Its downward slope means that the firms' outputs are
\keyterm{strategic substitutes}: when one firm produces more, the other
firm's best response is to produce less. In the example, every extra unit
chosen by the rival reduces firm $i$'s best-response quantity by one half.
The fixed cost of \euro10 affects profit, but not this conditional output
choice.

\section{Cournot--Nash equilibrium}

A \keyterm{Cournot--Nash equilibrium} is a pair of quantities such that each
firm's quantity is a best response to the other's quantity. It is a stable
prediction in the limited sense that neither firm wishes to change its own
quantity while taking the other's quantity as given. For the running example,
the two reaction functions intersect at the mutual best response shown in
Figure~\ref{fig:cournot-reaction-functions}.

\begin{figure}[htbp]
  \centering
  \begin{tikzpicture}[x=0.255cm,y=0.140cm]
    \draw[->,PPEInk] (0,0) -- (22.4,0) node[below,font=\small] {Firm 1 output $q_1$};
    \draw[->,PPEInk] (0,0) -- (0,40.8) node[left,font=\small] {Firm 2 output $q_2$};
    \draw[PPEBurgundy,very thick] (0,38) -- (19,0);
    \draw[PPENavarraRed,very thick] (0,19) -- (22,8);
    \draw[PPEGold,densely dashed] (12.67,0) -- (12.67,12.67) -- (0,12.67);
    \fill[PPEGold] (12.67,12.67) circle (2pt);
    \node[font=\scriptsize,text=PPEBurgundy,fill=white,inner sep=1pt] at (15.7,7.2) {$BR_1$};
    \node[font=\scriptsize,text=PPENavarraRed,fill=white,inner sep=1pt] at (16.2,12.9) {$BR_2$};
    \node[below,font=\scriptsize,text=PPEGold] at (12.67,0) {$38/3$};
    \node[left,font=\scriptsize,text=PPEGold] at (0,12.67) {$38/3$};
    \node[font=\scriptsize,text=PPEGold,align=left] at (15.9,17.7) {Cournot--Nash\\equilibrium};
  \end{tikzpicture}
  \caption{Reaction functions for the two-firm Cournot model. Their
  intersection is the output pair at which each firm is choosing a best
  response to the other.}
  \label{fig:cournot-reaction-functions}
\end{figure}
\FloatBarrier

With identical firms, set $q_1=q_2=q$ in the reaction function:
\[
q=19-\frac{1}{2}q,
\]
so
\[
q_1^*=q_2^*=\frac{38}{3}.
\]
Total output and price are
\[
Q^*=\frac{76}{3},
\qquad
p^*=\frac{88}{3}.
\]
Each firm's equilibrium profit is
\[
\pi_i^*=\frac{2798}{9}.
\]

At equilibrium, neither firm can increase profit by changing its output alone.
This does not mean that firms cooperate or that the outcome maximises joint
profit. Each firm ignores the effect of its output decision on the rival's
profit. A Cournot--Nash equilibrium is a mutual-best-response outcome; the
later game-theory chapter develops this idea for a wider range of strategic
settings.

\section{Comparison with monopoly and competition}

In a monopoly, a single firm restricts output to maximise industry profit. In
Cournot oligopoly, each additional firm expands total output and pushes price
toward marginal cost. The model therefore gives a useful middle benchmark:
more competitive than monopoly, but not identical to perfect competition.
With the running demand curve $p=80-2Q$ and $MC=4$, the three static
benchmarks are shown in Figure~\ref{fig:market-structure-comparison}:
\[
Q_M<Q_{\text{Cournot}}<Q_{PC},
\qquad
p_M>p_{\text{Cournot}}>p_{PC}.
\]

\begin{figure}[htbp]
  \centering
  \begin{tikzpicture}[x=0.155cm,y=0.065cm]
    \draw[->,PPEInk] (0,0) -- (42.5,0) node[below,font=\small] {Total output $Q$};
    \draw[->,PPEInk] (0,0) -- (0,83) node[left,font=\small] {Price};
    \draw[PPENavarraRed,very thick,domain=0:40,samples=80]
      plot (\x,{80-2*\x});
    \draw[PPEBurgundy,very thick] (0,4) -- (40.5,4);
    \draw[PPEGold,densely dashed] (0,42) -- (19,42) -- (19,0);
    \draw[PPEGold,densely dashed] (0,29.33) -- (25.33,29.33) -- (25.33,0);
    \draw[PPEGold,densely dashed] (38,0) -- (38,4);
    \fill[PPEGold] (19,42) circle (2pt);
    \fill[PPEGold] (25.33,29.33) circle (2pt);
    \fill[PPEGold] (38,4) circle (2pt);
    \node[left,font=\scriptsize,text=PPEGold] at (0,42) {$p_M$};
    \node[left,font=\scriptsize,text=PPEGold] at (0,29.33) {$p_{\mathit{Cournot}}$};
    \node[left,font=\scriptsize,text=PPEGold] at (0,4) {$p_C=MC$};
    \node[below,font=\scriptsize,text=PPEGold] at (19,0) {$Q_M$};
    \node[below,font=\scriptsize,text=PPEGold] at (25.33,0) {$Q_{\mathit{Cournot}}$};
    \node[below,font=\scriptsize,text=PPEGold] at (38,0) {$Q_C$};
    \node[font=\scriptsize,text=PPENavarraRed,fill=white,inner sep=1pt] at (31.8,18.2) {Demand};
    \node[font=\scriptsize,text=PPEBurgundy,fill=white,inner sep=1pt] at (28.2,8.0) {$MC$};
    \node[font=\scriptsize,text=PPEBurgundy] at (17.2,50.0) {Monopoly};
    \node[font=\scriptsize,text=PPEBurgundy] at (27.7,35.6) {Cournot};
    \node[font=\scriptsize,text=PPEBurgundy,fill=white,inner sep=1pt,align=right] at (42.0,15.0) {Perfect\\competition};
  \end{tikzpicture}
  \caption{Monopoly, Cournot duopoly, and perfect competition. Cournot
  output and price lie between the monopoly and competitive benchmarks.}
  \label{fig:market-structure-comparison}
\end{figure}
\FloatBarrier

With $N$ identical active firms, inverse demand $p=a-bQ$, constant marginal
cost $c$, and $a>c$, the symmetric Cournot solution is
\[
q_i^*=\frac{a-c}{b(N+1)},
\qquad
Q^*=\frac{N(a-c)}{b(N+1)},
\qquad
p^*=\frac{a+Nc}{N+1}.
\]
As $N$ rises, output approaches the competitive level $(a-c)/b$ and price
approaches $c$. This is a result within the model, not a promise that every
real concentrated market behaves this way. Fixed costs do not alter these
conditional output choices, but they can affect whether entry is profitable
and thus the number of active firms.

The number of firms alone does not determine market performance. Entry
barriers, capacity, product differentiation, repeated interaction, and the
possibility of collusion also matter.

\section{Collusion}

Firms can increase joint profit by restricting output as a monopoly would. An
agreement to do so is \keyterm{collusion}; a group of firms that coordinates
in this way is often called a \keyterm{cartel}. Each firm would prefer the
high joint price if all comply, yet it may gain even more by secretly selling
extra units while rivals restrict their output. Collusion is therefore
difficult to sustain because every member has an incentive to secretly expand
output or cut price.

\begin{center}
\small
\begin{tabular}{c c c}
\toprule
 & \multicolumn{2}{c}{Firm 2} \\
Firm 1 & Restrict output & Expand output \\
\midrule
Restrict output & $(\pi_C,\pi_C)$ & $(\pi_L,\pi_D)$ \\
Expand output & $(\pi_D,\pi_L)$ & $(\pi_N,\pi_N)$ \\
\bottomrule
\end{tabular}

\vspace{0.35em}
$\pi_D>\pi_C>\pi_N>\pi_L$: secretly expanding is tempting when the rival
restricts output.
\end{center}

Competition policy commonly prohibits agreements that reduce competition,
although assessing conduct in real markets can be complex. The payoff logic
above foreshadows the prisoner's-dilemma structure analysed in Chapter~17.

\section*{Chapter Review}

\subsection*{Key concepts and equations}
\begin{itemize}
  \item \textbf{Oligopoly}: a market with few strategically interdependent firms.
  \item \textbf{Reaction function}: a firm's best response for each rival action.
  \item \textbf{Cournot competition}: firms choose quantities simultaneously.
  \item \textbf{Cournot--Nash equilibrium}: each quantity is a best response to the other.
  \item \textbf{Strategic substitutes}: more output by one firm reduces the other's best-response output.
  \item \textbf{Collusion}: coordination intended to raise joint profits by limiting competition.
\end{itemize}

\section*{Exercises}

\subsection*{Conceptual questions}
\begin{enumerate}[label=\arabic*.]
  \item What does strategic interdependence mean? Give two examples outside quantity competition.
  \item Explain why treating a rival's output as fixed in a Cournot choice does not mean believing that the rival will never react.
  \item When is Cournot a more useful model than Bertrand competition?
  \item Why might a cartel be unstable even if all firms gain from cooperation compared with competition?
\end{enumerate}

\subsection*{Graphical exercises}
\begin{enumerate}[resume,label=\arabic*.]
  \item Draw two downward-sloping reaction functions. Mark their intersection and explain its meaning.
  \item On a demand diagram, compare monopoly, Cournot duopoly, and perfect competition. Label the three quantities and prices.
\end{enumerate}

\subsection*{Analytical exercises}
\begin{enumerate}[resume,label=\arabic*.]
  \item Two firms face $p=80-2Q$ and $C(q_i)=10+4q_i$. Derive both reaction functions and solve the Cournot equilibrium.
  \item $N$ identical firms face $p=120-3Q$ and constant marginal cost $10$. Find total output and price as functions of $N$. Evaluate the monopoly case, $N=5$, and the limit as $N\to\infty$.
  \item Explain how a fall in marginal cost for only one Cournot firm affects its reaction function and the equilibrium.
\end{enumerate}

\subsection*{Exam-style multiple-choice questions}
\begin{enumerate}[resume,label=\arabic*.]
  \item In the standard Cournot model, firms choose: (a) prices; (b) output quantities; (c) wages; or (d) consumer preferences.
  \item A Nash equilibrium occurs when: (a) firms maximise joint profit; (b) no firm benefits from a unilateral change; (c) firms charge marginal cost; or (d) all firms are identical.
  \item As the number of identical Cournot firms rises, price tends to: (a) rise toward the monopoly price; (b) approach marginal cost; (c) become unrelated to cost; or (d) become zero regardless of demand.
\end{enumerate}

\subsection*{PPE discussion question}

\noindent Is a concentrated market necessarily harmful to society? Discuss how
the answer depends on innovation, economies of scale, product quality, and the
ability of institutions to prevent collusion or abuse of market power.
